% Hyperellipticity of a Shimura curve quotient from several Borcherds divisors, by counting against
% the rational torsion of the Jacobian, with no appeal to Guo-Yang's Assumption 40.
%
% The argument is elementary once the two inputs are in place -- Borcherds' finite-order character
% and a bound on the rational torsion -- and the point of writing it down is that both inputs are
% computable here without any model of the curve, so the criterion is effective.
%
% Status: the mathematics below is complete as stated. Its application to X_0^46(31)* rests on a
% search whose exhaustiveness is not part of the argument (see Remark rem:exhaust), and on the
% banked certificates described in Section sec:effective.
%
% Numbering is local to this file.
\documentclass[11pt,reqno]{amsart}
\usepackage[margin=1.15in]{geometry}
\usepackage{amsmath,amssymb,amsthm,mathtools}
\usepackage{booktabs}
\usepackage[colorlinks=true,linkcolor=blue,citecolor=blue,urlcolor=blue]{hyperref}
\theoremstyle{plain}
\newtheorem{theorem}{Theorem}
\newtheorem{proposition}[theorem]{Proposition}
\newtheorem{lemma}[theorem]{Lemma}
\newtheorem{corollary}[theorem]{Corollary}
\theoremstyle{definition}
\newtheorem{remark}[theorem]{Remark}
\newtheorem{example}[theorem]{Example}
\DeclareMathOperator{\Pic}{Pic}
\DeclareMathOperator{\Jac}{Jac}
\DeclareMathOperator{\divisor}{div}
\newcommand{\Q}{\mathbb{Q}}
\newcommand{\Z}{\mathbb{Z}}
\newcommand{\PP}{\mathbb{P}}
\begin{document}
\title{Hyperellipticity from several Borcherds divisors, by counting torsion}
\author{}
\date{\today}
\maketitle

\section{The problem this removes}

Guo and Yang's method for the equation of a Shimura curve quotient produces functions as Borcherds
lifts: a weakly holomorphic form $f$ of weight $1/2$ with prescribed principal part lifts to a
meromorphic modular form $\Psi_f$ whose divisor is the Heegner divisor $Z(f)$ prescribed with it.
What the lift does not give for free is that $\Psi_f$ is a \emph{function} on the quotient. Their
Assumption~40 asks exactly this, and it is a genuine hypothesis: the character of $\Psi_f$ has
finite order, so a power of $\Psi_f$ is a function and the class of $Z(f)$ is torsion in the
Jacobian, but the class need not be trivial.

So a search that returns one divisor $Z(f)$ of degree two proves hyperellipticity only under the
assumption. The observation below is that a search returning \emph{several} divisors proves it
outright, as soon as there are more of them than the rational torsion of the Jacobian can separate.
The hypothesis is then replaced by a count, and both sides of the count are computable without a
model of the curve.

\section{The criterion}

Throughout, $X$ is a smooth projective geometrically connected curve over $\Q$ of genus $g \ge 2$
with $X(\Q) \neq \emptyset$, and $J = \Jac(X)$. Since $X$ has a rational point, the natural map
$\Pic^0(X)(\Q) \to J(\Q)$ is an isomorphism, so a $\Q$-rational divisor class of degree zero is a
point of $J(\Q)$.

\begin{lemma}[Borcherds]\label{lem:torsion}
Let $f$ be a weakly holomorphic form whose principal part is integral, and let $\Psi_f$ be its
Borcherds lift on $X$, with divisor $Z(f)$. Then the class of $Z(f)$ is torsion in $J(\Q)$.
\end{lemma}

\begin{proof}
Integrality of the principal part is the hypothesis under which the character of $\Psi_f$ has
finite order, say $n$; see \cite[proof of Theorem 4.1]{borcherds-gkz}, with the correction
\cite{borcherds-correction} restoring the finite-order condition that the printed statement omits
from the definition of a principal Heegner divisor. Then $\Psi_f^{\,n}$ has trivial character and
is a rational function on $X$, so $n \cdot Z(f) = \divisor(\Psi_f^{\,n})$ is principal. The divisor
$Z(f)$ is a sum of Heegner cycles and so is Galois-stable, hence defines a class in
$\Pic^0(X)(\Q) = J(\Q)$, which is therefore killed by $n$.
\end{proof}

\begin{corollary}\label{cor:rational}
Let $f$ have \emph{rational} principal part, with $n$ a positive integer clearing its
denominators, and suppose the Borcherds lift of $nf$ has divisor $n Z$. Then the class of $Z$ is
torsion in $J(\Q)$.
\end{corollary}

\begin{proof}
$nf$ has integral principal part, so by Lemma~\ref{lem:torsion} the class of $nZ$ is torsion. A
class whose $n$-th multiple is torsion is itself torsion, $J(\Q)$ being finitely generated.
\end{proof}

\begin{theorem}\label{thm:count}
Let $t$ be an integer with $\#J(\Q)_{\mathrm{tors}} \le t$. Suppose $D_1, \dots, D_r$ are pairwise
distinct effective divisors of degree two on $X$, each defined over $\Q$, such that for all $i \ne
j$ some positive integer multiple of $D_i - D_j$ is the divisor of a Borcherds form with integral
principal part. If $r > t$, then $X$ is hyperelliptic.
\end{theorem}

\begin{proof}
By Lemma~\ref{lem:torsion} and Corollary~\ref{cor:rational} the classes $[D_i - D_1]$, for
$i = 1, \dots, r$, all lie in
$J(\Q)_{\mathrm{tors}}$, a set of at most $t$ elements. As $r > t$, two of them coincide: there are
$i \ne j$ with $[D_i - D_1] = [D_j - D_1]$, hence $[D_i - D_j] = 0$ and $D_i \sim D_j$.

Since $D_i \ne D_j$, the linear equivalence is realised by a nonconstant $h \in \Q(X)$ with
$\divisor(h) = D_i - D_j$. Write $G$ for the greatest common divisor of the effective divisors $D_i$
and $D_j$, so that the polar divisor of $h$ is $D_j - G$, of degree $2 - \deg G$ with $\deg G \in
\{0, 1\}$; it is not $2$, as $D_i \neq D_j$. The degree of $h$ as a morphism $X \to \PP^1$ is the
degree of its polar divisor. If $\deg G = 1$ that degree is one and $X \cong \PP^1$, contradicting
$g \ge 2$. Hence $\deg G = 0$ and $h$ has degree two, so $X$ is hyperelliptic.
\end{proof}

\begin{remark}
Only $r > t$ is used, not which pair of divisors becomes equivalent, and the proof does not
identify the pair. For the \emph{equation} of $X$ one must know it, and there the candidates are
few and can be separated by a different means; see Remark~\ref{rem:which}.
\end{remark}

\begin{remark}
The argument is not special to degree two. With $D_1, \dots, D_r$ effective of degree $d$ and
$r > t$, the same count gives $D_i \sim D_j$ and hence a map $X \to \PP^1$ of degree $d - \deg G
\le d$, so the gonality of $X$ is at most $d$. Degree two is the case where the conclusion is a
dichotomy, because a map of degree one is excluded by the genus.
\end{remark}

\section{Why both sides of the count are computable}\label{sec:effective}

The criterion is useful here because neither hypothesis needs a model of $X$.

\emph{The torsion bound.} For an abelian variety over $\Q$ with good reduction at an odd prime $p$,
reduction is injective on the rational torsion \cite[Appendix]{katz-torsion}. Hence
$\#J(\Q)_{\mathrm{tors}}$ divides $\#J(\mathbb{F}_p) = P_p(1)$, where $P_p$ is the numerator of the
zeta function of $X$ over $\mathbb{F}_p$, and so divides the greatest common divisor of these over
any set of such primes. For a quotient of a Shimura curve the point counts over $\mathbb{F}_{p^k}$
come from the trace formula, so $P_p$ and the bound are computed with no equation for $X$ in sight.

\emph{The divisors.} Each $D_i - D_j$ is produced by a search for a weakly holomorphic form with
that prescribed Heegner divisor. Two things are then to be checked, and both are checked on the
form itself rather than on the search: the denominator $n$ of its principal part, after which
Corollary~\ref{cor:rational} applies --- in practice the search returns a rational principal part,
its integrality option making the solution integral over an echelonised basis of the pool, which is
a weaker thing --- and that a form with that principal part exists at all. The second is
witnessed by presenting the form as a combination of eta quotients, which is what the search
computes and what the existence claim really rests on. The coefficients of that combination are
rational, and in practice large --- they come out of an echelon form --- so the presentation is not
short; but it is finite and explicit, and expanding it and recomputing the divisor is cheap
compared with the search that produced it. Recording it is therefore what lets a later reader
verify the hypothesis without repeating the search.

\begin{remark}\label{rem:exhaust}
The count needs $r$ divisors to exist, not that the search found all there are. Exhaustiveness of a
search is therefore never part of the argument, which matters because exhaustiveness is the part
that is expensive and the part that depends on a chosen window of pole orders.
\end{remark}

\section{The case that prompted this}

\begin{example}
Let $X = X_0^{46}(31)/W_{46,31}$, of genus three, which has no Atkin--Lehner involution left and so
is reached by none of the cover constructions. A search over the degree-two divisors supported on
its rational CM points and its quadratic CM cycles returns exactly six, and they are the six
pairwise differences of
\[
  A = P_{-3} + P_{-403}, \quad B = P_{-24} + P_{-232}, \quad C = P_{-115} + P_{-372}, \quad
  D = Z(-292),
\]
the last a cycle of degree two; the six principal parts are rational with denominators $20, 20,
10, 10, 4, 10$, so Corollary~\ref{cor:rational} applies to each. The trace
formula gives $\#J(\mathbb{F}_p)$ equal to $86, 356, 604, 1882, 3554, 5560$ at
$p = 3, 5, 7, 11, 13, 17$, whose greatest common divisor is $2$, so $t = 2$. Four effective
divisors against $t = 2$ satisfies $r > t$ with room to spare, and Theorem~\ref{thm:count} gives
that $X$ is hyperelliptic, unconditionally.
\end{example}

\begin{remark}\label{rem:which}
Which of the six differences is principal is not determined by the count. For the equation one
takes each candidate in turn as the degree-two coordinate, builds the model it gives, and compares
the number of points of that model over $\mathbb{F}_p$ with the trace formula, which again needs no
prior model and which a wrong candidate fails.
\end{remark}

\begin{remark}
The six divisors have pole orders at infinity between $556$ and $1760$, and three of them use the
quadratic cycle. Neither is reachable by a search restricted to divisors with a single double pole,
which is how such searches were previously posed; that restriction returns nothing at all here. It
is the availability of several divisors, not of one deeper one, that removes the assumption.
\end{remark}

\begin{thebibliography}{9}
\bibitem{borcherds-gkz} R. E. Borcherds, \emph{The Gross--Kohnen--Zagier theorem in higher
dimensions}, Duke Math. J. \textbf{97} (1999), 219--233.
\bibitem{borcherds-correction} R. E. Borcherds, \emph{Correction to: ``The Gross--Kohnen--Zagier
theorem in higher dimensions''}, Duke Math. J. \textbf{105} (2000), 183--184.
\bibitem{guo-yang} J.-W. Guo and Y. Yang, \emph{Equations of hyperelliptic Shimura curves},
Compositio Math. \textbf{153} (2017), 1--40.
\bibitem{katz-torsion} N. M. Katz, \emph{Galois properties of torsion points on abelian varieties},
Invent. Math. \textbf{62} (1981), 481--502.
\end{thebibliography}
\end{document}
